% Extract: the summary points of the fact check of the CBC Daybreak
% interview of October 9, 2026, verbatim (filled by build.py).
\documentclass[10pt,letterpaper]{article}
\usepackage[T1]{fontenc}
\usepackage[utf8]{inputenc}
\usepackage[scaled=0.95]{helvet}
\renewcommand{\familydefault}{\sfdefault}
\usepackage[letterpaper,top=0.8in,bottom=0.8in,left=0.9in,right=0.9in]{geometry}
\usepackage{microtype}
\usepackage{enumitem}
\usepackage{titlesec}
\usepackage{fancyhdr}
\usepackage{xcolor}
\usepackage{booktabs}
\usepackage{longtable}
\usepackage{array}
\usepackage{textcomp}
\usepackage[hidelinks]{hyperref}
\urlstyle{same}
\newcommand{\docversion}{2026-10-16-15-38-58} % set by build.py (America/Toronto)
\hypersetup{pdfsubject={version \docversion}}

\setlength{\parindent}{0pt}
\setlength{\parskip}{0.5em}
\linespread{1.0}

\titleformat{\section}{\normalfont\large\bfseries}{\thesection}{0.6em}{}[\vspace{1pt}\titlerule]
\titlespacing*{\section}{0pt}{1.1em}{0.4em}
\titleformat{\subsection}{\normalfont\bfseries}{\thesubsection}{0.6em}{}
\titlespacing*{\subsection}{0pt}{0.8em}{0.25em}

\pagestyle{fancy}
\fancyhf{}
\renewcommand{\headrulewidth}{0pt}
\fancyfoot[L]{\footnotesize\color{gray}CBC Daybreak, October 9, 2026 -- fact check, summary points -- \repourl}
\fancyfoot[R]{\tiny\color{gray}v.\,\docversion}

% ---- Links ----
\newcommand{\repourl}{https://github.com/OpenMobilityData/Terrebonne-open-sources}
\newcommand{\repo}[1]{\href{\repourl}{#1}}
\newcommand{\doc}[2]{\href{\repourl/blob/main/#1}{#2}}
\newcommand{\arurl}{https://github.com/OpenMobilityData/Terrebonne-access-requests}
\newcommand{\pagesurl}{https://openmobilitydata.github.io/Terrebonne-access-requests}
\newcommand{\cliptime}[1]{\textsf{\small[#1]}}
% \pg{pages}: page of bundle 01 of access request 2026-581, linked
\newcommand{\pgl}[2]{\href{\pagesurl/2026-581/raw/01_emails_mandate_counts.pdf\#page=#1}{\textsf{\small[2026-581, p.~#2]}}}
\newcommand{\pg}[1]{\pgl{#1}{#1}}
% \exp{page}: page of the 2023 EXP study (open-sources repository)
\newcommand{\exppg}[1]{\href{\repourl/blob/main/2023-exp-mobility-study/raw/EXP_Etude_de_mobilite_rue_de_Terrebonne_2023.pdf}{\textsf{\small[EXP 2023, p.~#1]}}}
\newcommand{\quot}[1]{\begin{quote}\small #1\end{quote}}

\begin{document}

{\large\bfseries CBC \emph{Daybreak}, October 9, 2026: the councillor's statements checked against the record}\hfill{\small October~16, 2026}

\vspace{-0.2em}
{\small Summary points of the fact check, verbatim. Full document, sources, page references and the transcript:\\ \href{\repourl/tree/main/2026-cbc-daybreak-interview}{\repourl/tree/main/2026-cbc-daybreak-interview}.}
\vspace{0.4em}

\begin{enumerate}[nosep,itemsep=0.3em,leftmargin=1.4em]
  \item \begin{minipage}[t]{\linewidth}\textbf{``Less than one an hour'' in winter
    \cliptime{03:11}: partly consistent.} Over December~1 -- March~31 the
    borough's counter at Kensington recorded 4{,}891 bicycle passages in 121
    days: 40.4 a day, or 1.68 an hour counted around the clock, both
    directions. Taken one direction at a time over all 24 hours the rate is
    below one an hour (eastbound 0.96, westbound 0.73, the westbound channel
    undercounting from about March~1). On weekdays between 07:00 and 19:00 it
    was 3.2 an hour. On the City's open data, the one Terrebonne channel
    (westbound, at Girouard) averaged 27.4 a day, 1.14 an hour, over the same
    months.\end{minipage}\vspace{0.35em}
  \item \begin{minipage}[t]{\linewidth}\textbf{The ``8'' was ``the sum of all the hours during the rush
    hour… six hours'' \cliptime{07:03}: consistent.} The figure is on the
    borough's slide \pg{62}: bicycles on the Terrebonne axis at one
    intersection, Grand, on one day, January~29, 2025, over the six peak
    hours defined at \pg{34} (06:30--09:30 and 15:30--18:30), both
    directions together. The same slide gives 10 at Cavendish that day. In
    his October~6 address the figure was described as ``8 passage, right? So
    that's 4 cyclists.''\end{minipage}\vspace{0.35em}
  \item \begin{minipage}[t]{\linewidth}\textbf{The open data website is ``connected directly to
    Eco-Counter, which has those loops installed on the ground''
    \cliptime{07:19}: not consistent with the open data.} The City's
    \emph{Compteurs cyclistes permanents} dataset lists 128 counting sites,
    77 of them Eco-Counter; none is on Terrebonne. The only Terrebonne site
    is a City detector at Girouard (source ``Ville''), whose one Terrebonne
    channel counts westbound traffic, with daily data from October~23, 2025.
    The borough's Eco-Counter loop at Kensington is not on the portal; its
    data have been made public only through access requests 2025-697,
    2026-300 and 2026-581.\end{minipage}\vspace{0.35em}
  \item \begin{minipage}[t]{\linewidth}\textbf{``From 40 a day to 30 a day during December to March''
    \cliptime{07:34}: partly consistent.} The figures match the borough's
    Kensington counter, both directions: December 40.0, January 32.7,
    February 34.9 a day; March was 53.5. On the open data channel named as
    the source (westbound only, at Girouard): December 32.5, January 22.6,
    February 29.1, March 24.7.\end{minipage}\vspace{0.35em}
  \item \begin{minipage}[t]{\linewidth}\textbf{``We have the data for many years'' \cliptime{03:33}:
    partly consistent.} Continuous counts on Terrebonne cover one winter,
    2025--26 (Kensington loop from July~26, 2025; open-data detector from
    October~23, 2025). Earlier winters are covered by single-day intersection
    counts (Miovision, November 2024 and January 29, 2025) and by Telraam
    counts ending in December 2024 \pg{65}.\end{minipage}\vspace{0.35em}
  \item \begin{minipage}[t]{\linewidth}\textbf{``If you go back to 2023… they basically removed about 250
    parking spots'' \cliptime{01:08}: partly consistent.} The study and the
    announcement date from 2023; the lane was approved in February 2024 and
    built in May--June 2024. The 2023 EXP study counts 475 spaces on
    Terrebonne, 375 once the 5-metre clearances are applied, and 171 in the
    retained design \exppg{100}: a reduction of 304, or 204 against the
    375. No as-built count is in the record. In 2020 a press report gave 250
    spaces removed by the temporary lanes of that year (\emph{Métro},
    September~1, 2020).\end{minipage}\vspace{0.35em}
  \item \begin{minipage}[t]{\linewidth}\textbf{``At least 50 to 60 parking spots'' \cliptime{01:28}: not
    found in the record.} The figure does not appear in the 508 pages of
    request 2026-581, which contain no inventory of on-street spaces, no
    occupancy count and no estimate of the spaces the change would create.
    The only occupancy survey on record is EXP's, taken on two days in
    autumn 2022, before the lane, in the evening: 35--70\% on weekdays
    \exppg{54}.\end{minipage}\vspace{0.35em}
  \item \begin{minipage}[t]{\linewidth}\textbf{``We vote here at the borough, then a vote at the Conseil
    Municipal, then it comes back to the borough'' \cliptime{08:19}:
    consistent with the 2024 precedent.} Terrebonne is on the City's
    arterial network. In 2024 the borough offered the work under art.~85 of
    the City's Charter (CA24 170018, February~5), City council accepted it
    (CM24 0179, February~19), and the borough then adopted the one-way
    ordinance (CA24 170131, June~3): three votes over four months. No
    borough or City council agenda through October~19, 2026 carries a
    decision item on Terrebonne. Implementation ``before Christmas'' cannot
    be verified.\end{minipage}\vspace{0.35em}
  \item \begin{minipage}[t]{\linewidth}\textbf{``So for now, we're looking at it in a way that will make
    it permanent'' \cliptime{12:22}: differs from earlier descriptions.} On
    October~6 the councillor said: ``It's going to be a test.'' The
    borough's service request of August~20 describes the winter measure as
    ``mesure transitoire pour l'hiver 2026-2027'' \pg{4}; the permanent
    mandate of August~13 concerns a different layout, a two-way lane on one
    side \pg{1}.\end{minipage}\vspace{0.35em}
  \item \begin{minipage}[t]{\linewidth}\textbf{``The final study… still need to be submitted to us''
    \cliptime{00:38}: cannot be verified.} EXP's final note was due on
    October~2, 2026 \pg{12}; a draft was transmitted on September~17
    \pg{167}. No final note is in the release, and the engineering study and
    the City's analysis are withheld. A request for the final note is
    pending (answer due November~7).\end{minipage}\vspace{0.35em}
  \item \begin{minipage}[t]{\linewidth}\textbf{``We've been able to reduce vehicle traffic by two-thirds''
    \cliptime{09:56}: consistent with the range in the record.} Before the
    project: 6{,}300 vehicles a day (2022) between Cavendish and Grand
    \exppg{41}. After: 1{,}375 a day on the Kensington counter's first year,
    a reduction of 78\%; peak-hour entries at Grand fell by 46--63\%
    \pg{34}. The one-way change at Cavendish was part of the 2024
    project.\end{minipage}\vspace{0.35em}
  \item \begin{minipage}[t]{\linewidth}\textbf{Compared with October~6.} In his address to the borough
    council three days earlier, the councillor described the data as ``the
    official data provided by the borough'' (on \emph{Daybreak}: ``the open
    data website''); the winter change as a 99\,\% drop, ``from 700 during the
    good months to 8'' (on \emph{Daybreak}: 30 to 40 a day, and the 8 as a
    six-hour sum); ``data for 3 years in a row'' (``many years''); ``more than 30
    meetings'' (``more than 40''); and the measure as ``a test'' (``for now…
    permanent''). No number of parking spaces was given in the address. The
    statements are set side by side in section~3.\end{minipage}\vspace{0.35em}
\end{enumerate}

\end{document}
